\section{LIMITATIONS}
\label{app:limitations}
\readorder{3}

The conclusion names the main limitations. Each is argued where it arises in the paper; this appendix collects them all, with what each one implies.
\begin{itemize}
\item \emph{The form of the leakage.} The leak fraction is one number per section, the leakage is the same for every gene and scales with the receiving cell's depth, and there is no ambient term (\cref{sec:generative}). A per-cell or per-gene variation in contamination is not covered by the sweep; its sensitivity to the leak's form is checked on the primary section only, where three other forms move no read by more than $1.3$ seed standard deviations, but a fixed false-positive floor lowers the transport read from $0.65$ to $0.50$, a move of $8.3$ standard deviations (\cref{app:sensitivity}).
\item \emph{What the sweep separates.} It separates a spatial association from leakage of the modelled form, not from spatially organised intrinsic variation within a type, which the conditional invariance assigns to the response, nor from cells selecting their niche (\cref{sec:sweep}).
\item \emph{Labels.} The labels come from the same contaminated counts, and the decomposition is relative to their granularity: a state merged into its lineage becomes variation the model must place, and a spatially clustered state merged in this way is read as a response (\cref{sec:setting,sec:invariance}).
\item \emph{The invariance is incomplete.} With the adversary, a nonlinear probe still recovers $3.5$ to $5.5\%$ of the within-type variance of neighbour composition from the intrinsic state on every section; where little leaked to begin with, as on the TMA core, the adversary removes only $29$ to $48\%$ of what the uncontrolled model leaves. A probe at its floor bounds only the dependence it can detect (\cref{tab:probe}).
\item \emph{The response at the operating point.} With the response weight far above its bound-equivalent value, the response readouts are in effect readouts of the context regression, except that on the lung section the response separates some merged clusters better than its prior mean (\cref{app:subtype}); the per-cell channel is closed by design, not by the data: a per-cell response planted in simulation is not detected at this weight and is partly absorbed by the intrinsic state (\cref{sec:objective}, \cref{app:planted-percell}). Several response readouts are also not specific to it: ligand--receptor genes take a larger share of the leakage channel than other genes on every section, so that criterion is met by leakage alone; the lean of the programmes toward membrane and secreted genes is what a gene-specific leak would produce, and the band ordering on the primary section is also followed by the intrinsic state and by depth (\cref{app:response}).
\item \emph{The cycle reference.} The cell-cycle scores against which the intrinsic state is graded are gene-set scores of the cell's own segmented counts, not an independent measurement: transcripts leaked from cycling neighbours and expression induced by the environment enter them, so the asymmetry between the latents shows where the model places this axis of the counts, not that it is the cell's own (\cref{sec:results-z}).
\item \emph{Inference.} The intrinsic encoder does not see the foreign influx, so its posterior mean can differ from the exact one; in simulation, at the final configuration, the difference is measurable, $0.65$ to $1.93$ posterior standard deviations, and larger at every depth than that of a perceptron fitted to the exact posterior mean from the same inputs, and on tissue it cannot be measured (\cref{app:planted}). The posterior family is a compromise, so the uncertainty reported comes from the sweep and the seeds, not from the posterior (\cref{sec:inference}).
\item \emph{Identification.} Beyond the upper bound on the leak fraction and the invariance the probe grades, nothing guarantees that the intrinsic and response coordinates are identified (\cref{sec:sweep}).
\item \emph{Evaluation.} Cells of held-out tiles enter training as neighbours under a stop-gradient (\cref{sec:batching}), every fit is to one section, and only one section is held out whole. Most cells scored by the transport reads lie in the model's training tiles; scored on held-out tiles only, the fraction of the ceiling moves in both directions and Read A falls (\cref{tab:transport-heldout}). The comparison methods were fitted on our held-out cells too, as their software is designed, and at their published defaults, while DISCELL's weights were calibrated on these sections (\cref{sec:setup}); two of them could not be run on some of the sections on the resources currently available (\cref{tab:baselines}).
\end{itemize}
